\appendix
\section*{Appendix R6: Fixed-Point Theorem for the Interpretation Operator}

\subsection*{R6.0 Overview}

This appendix establishes the fixed-point theorem for the NDH
interpretation operator
\[
I_p : X_{42} \to \{c\},
\]
showing that $I_p$ is the unique transformation-invariant fixed point
under the full NDH operator algebra. The proof relies on the smooth
structure of the 42D interpretation manifold $X_{42}$ and the bundle
structure defined in Appendix R5.

\subsection*{R6.1 Preliminaries}

Let $X_{42}$ be the 42-dimensional smooth manifold defined in Appendix
R4.

Let $\mathcal{O}(X_{42})$ denote the class of admissible smooth
operators:
\[
\mathcal{O}(X_{42}) = \{ A : X_{42} \to X_{42} \mid A \text{ smooth} \}.
\]

Let $E$ be the interpretive-field bundle defined in Appendix R5, with
sections:
\[
\Gamma(E).
\]

Let the transformation stack be:
\[
T = 
\mathcal{S}_{\Xi} \circ
\mathcal{C}_{\Psi} \circ
\mathcal{E}_{\Phi} \circ
\mathcal{F}_{\Theta} \circ
\mathcal{R}_{\theta} \circ
I_{\mathrm{man}}.
\]

Let $\mathcal{I}_{\Omega} \in \Gamma(E)$ be an interpretive-field
section.

\subsection*{R6.2 Definition of Fixed Point}

An operator $F : X_{42} \to X_{42}$ has a fixed point $x^*$ if:
\[
F(x^*) = x^*.
\]

For the interpretation operator $I_p$, the fixed point is the constant
$c$:
\[
I_p(c) = c.
\]

\subsection*{R6.3 Transformation-Invariant Fixed Point}

We say that $I_p$ is transformation-invariant if:
\[
I_p(A(x)) = I_p(x)
\]
for all $A \in \mathcal{O}(X_{42})$.

\subsection*{R6.4 Theorem Statement}

\textbf{Theorem (42D Interpretation Fixed-Point Theorem).}  
Let $I_p : X_{42} \to \{c\}$ be the interpretation operator. Then:

\begin{enumerate}
    \item $I_p$ is a fixed point of the NDH operator algebra.
    \item $I_p$ is invariant under all admissible operators:
    \[
    I_p(A(x)) = I_p(x) \quad \forall A \in \mathcal{O}(X_{42}).
    \]
    \item $I_p$ is invariant under the full transformation stack:
    \[
    I_p(\mathcal{I}_{\Omega}(T(x))) = I_p(x).
    \]
    \item $I_p$ is the \emph{unique} operator satisfying (1)–(3).
\end{enumerate}

\subsection*{R6.5 Proof}

\paragraph{(1) Fixed Point}

Since $I_p$ is constant:
\[
I_p(x) = c,
\]
we have:
\[
I_p(c) = c,
\]
so $c$ is a fixed point.

\paragraph{(2) Invariance Under Admissible Operators}

Let $A \in \mathcal{O}(X_{42})$ be any smooth operator. Then:
\[
I_p(A(x)) = c = I_p(x).
\]
Thus $I_p$ is invariant under all admissible operators.

\paragraph{(3) Invariance Under the Transformation Stack}

Let $T$ be the transformation stack. Since $T : X_{42} \to X_{42}$ is
smooth:
\[
I_p(T(x)) = c = I_p(x).
\]

Since $\mathcal{I}_{\Omega}(T(x)) \in E_{T(x)}$ and $I_p$ induces a
trivial pullback bundle:
\[
I_p(\mathcal{I}_{\Omega}(T(x))) = c = I_p(x).
\]

\paragraph{(4) Uniqueness}

Suppose there exists another operator $J : X_{42} \to Y$ satisfying:
\[
J(A(x)) = J(x)
\]
for all $A \in \mathcal{O}(X_{42})$.

Let $A$ be the constant operator:
\[
A(x) = x_0.
\]

Then:
\[
J(x_0) = J(x).
\]

Since $x$ is arbitrary, $J$ must be constant:
\[
J(x) = y_0.
\]

Thus $I_p$ is the unique transformation-invariant fixed point.

\subsection*{R6.6 Consequences}

\begin{itemize}
    \item Interpretation is stable under all NDH transformations.
    \item Interpretive drift is impossible.
    \item Interpretive recursion collapses instantly.
    \item Interpretive fields cannot alter identity.
    \item The 42D layer is interpretively self-consistent.
\end{itemize}

\subsection*{R6.7 Summary}

The interpretation operator $I_p$ is the unique transformation-invariant
fixed point of the NDH operator algebra on the 42D interpretation
manifold. This establishes the mathematical foundation for 42D
interpretive stability.

This completes Appendix R6.
